%%%
%%% Radical "⼎"
%%%
\section*{Radical 15: ``⼎''}\addcontentsline{toc}{section}{Radical 15: ⼎}\addcontentsline{loh}{figure}{\#\#\#\# 15: ⼎}

%%%%%%%%%% 冬 %%%%%%%%%%
\subsection*{冬}\addcontentsline{loh}{figure}{冬}

\begin{Entry}{冬}{5}{⼎}
  \begin{Phonetics}{冬}{dong1}
    \definition*{s.}{Sobrenome: Dong}
    \definition{s.}{inverno}
    \definition{s.}{onomatopéia: som de um tambor batendo, batendo na porta, etc.}
  \end{Phonetics}
\end{Entry}

\begin{Entry}{冬天}{5,4}{⼎,⼤}
  \begin{Phonetics}{冬天}{dong1tian1}[][HSK 2]
    \definition[个]{s.}{inverno; a quarta estação do ano, na China, geralmente se refere aos três meses entre outubro e dezembro do calendário lunar}
  \seealsoref{冬}{dong1}
  \synonymref{冬季}{dong1ji4}
  \antonymref{夏季}{xia4ji4}
  \antonymref{夏天}{xia4tian1}
  \end{Phonetics}
\end{Entry}

\begin{Entry}{冬瓜}{5,5}{⼎,⽠}
  \begin{Phonetics}{冬瓜}{dong1gua1}
    \definition[个,斤,棵,株]{s.}{cabaça"-de"-cera; cabaça"-branca (fruto ou planta)}
  \end{Phonetics}
\end{Entry}

\begin{Entry}{冬季}{5,8}{⼎,⼦}
  \begin{Phonetics}{冬季}{dong1ji4}[][HSK 4]
    \definition[个,次,种]{s.}{inverno; o quarto trimestre do ano, habitualmente referido na China como o período de três meses entre o início do inverno e o início da primavera, e também referido como ``décimo, décimo primeiro e décimo segundo'' meses do calendário lunar}
  \synonymref{冬天}{dong1tian1}
  \antonymref{夏季}{xia4ji4}
  \end{Phonetics}
\end{Entry}

%%%%%%%%%% 冰 %%%%%%%%%%
\subsection*{冰}\addcontentsline{loh}{figure}{冰}

\begin{Entry}{冰}{6}{⼎}
  \begin{Phonetics}{冰}{bing1}[][HSK 4]
    \definition[块,层,些]{s.}{gelo; água em estado sólido |  algo parecido com gelo | (gíria) metanfetamina}
    \definition{v.}{colocar gelo; colocar gelo ao redor; colocar no gelo; resfriar objetos com gelo ou água fria | sentir frio}
  \synonymref{冷}{leng3}
  \synonymref{凉}{liang2}
  \antonymref{暖}{nuan3}
  \antonymref{热}{re4}
  \antonymref{炭}{tan4}
  \antonymref{温}{wen1}
  \end{Phonetics}
\end{Entry}

\begin{Entry}{冰山}{6,3}{⼎,⼭}
  \begin{Phonetics}{冰山}{bing1shan1}[][HSK 7-9]
    \definition[座]{s.}{montanha gelada; montanha coberta de gelo | \emph{iceberg}; enormes blocos de gelo flutuando no mar | Figurativo: indivíduo ou grupo em que não se pode confiar por muito tempo; uma metáfora para um poder em que não se pode confiar por muito tempo}
  \end{Phonetics}
\end{Entry}

\begin{Entry}{冰天雪地}{6,4,11,6}{⼎,⼤,⾬,⼟}
  \begin{Phonetics}{冰天雪地}{bing1tian1-xue3di4}
    \definition{expr.}{gelo e neve; um mundo de gelo e neve; descreve uma cena de neve e gelo cobrindo o céu, extremamente frio; clima extremamente frio}
  \antonymref{四季如春}{si4ji4-ru2chun1}
  \end{Phonetics}
\end{Entry}

\begin{Entry}{冰冷}{6,7}{⼎,⼎}
  \begin{Phonetics}{冰冷}{bing1leng3}
    \definition{adj.}{gelado; tão frio quanto gelo; muito frio | hostil; indiferente}
  \end{Phonetics}
\end{Entry}

\begin{Entry}{冰球}{6,11}{⼎,⽟}
  \begin{Phonetics}{冰球}{bing1qiu2}
    \definition[个]{s.}{hóquei no gelo | disco; a ``bola'' usada no hóquei no gelo}
  \end{Phonetics}
\end{Entry}

\begin{Entry}{冰雪}{6,11}{⼎,⾬}
  \begin{Phonetics}{冰雪}{bing1xue3}[][HSK 4]
    \definition{adj.}{puro como gelo e neve; descreve uma pessoa como pura}
    \definition[片,场]{s.}{gelo e neve}
  \synonymref{纯洁}{chun2jie2}
  \end{Phonetics}
\end{Entry}

\begin{Entry}{冰棍}{6,12}{⼎,⽊}
  \begin{Phonetics}{冰棍}{bing1gun4}
    \definition[根,种,支]{s.}{picolé}
  \end{Phonetics}
\end{Entry}

\begin{Entry}{冰棍儿}{6,12,2}{⼎,⽊,⼉}
  \begin{Phonetics}{冰棍儿}{bing1gun4r5}[][HSK 7-9]
    \definition[根,种,支]{s.}{picolé; pirulito congelado}
  \end{Phonetics}
\end{Entry}

\begin{Entry}{冰箱}{6,15}{⼎,⾋}
  \begin{Phonetics}{冰箱}{bing1xiang1}[][HSK 4]
    \definition[台,个]{s.}{geladeira; freezer; refrigerador; aparelhos para congelar alimentos ou medicamentos com gelo para mantê"-los frios}
  \end{Phonetics}
\end{Entry}

\begin{Entry}{冰激凌}{6,16,10}{⼎,⽔,⼎}
  \begin{Phonetics}{冰激凌}{bing1ji1ling2}
    \definition[份,盒,种,个]{s.}{sorvete}
  \end{Phonetics}
\end{Entry}

\begin{Entry}{冰糕}{6,16}{⼎,⽶}
  \begin{Phonetics}{冰糕}{bing1gao1}
    \definition[块,盒]{s.}{sorvete | picolé; geladinho}
  \seealsoref{冰棍儿}{bing1gun4r5}
  \synonymref{雪糕}{xue3gao1}
  \end{Phonetics}
\end{Entry}

%%%%%%%%%% 冲 %%%%%%%%%%
\subsection*{冲}\addcontentsline{loh}{figure}{冲}

\begin{Entry}{冲}{6}{⼎}
  \begin{Phonetics}{冲}{chong1}[][HSK 4]
    \definition{s.}{via pública; local importante; via de passagem; via local importante | um trecho de planície em uma área montanhosa | (astronomia) oposição; os planetas externos orbitam até ficarem alinhados com a Terra e o Sol, e a Terra está no meio}
    \definition{v.}{atacar; apressar; correr; passar rapidamente; passar por um obstáculo | colidir; chocar; bater | despejar água fervente sobre | enxaguar; dar descarga; lavar | revelar (filme) | neutralizar a má sorte}
  \synonymref{冲破}{chong1po4}
  \synonymref{冲突}{chong1tu1}
  \synonymref{冲洗}{chong1xi3}
  \synonymref{小}{xiao3}
  \synonymref{幼}{you4}
  \end{Phonetics}
  \begin{Phonetics}{冲}{chong4}[][HSK 6]
    \definition{adj.}{poderoso; com vigor; com muita força; vigoroso | forte; odor forte e pungente (olfato)}
    \definition{prep.}{de frente; em direção a | na força de; com base em; em virtude de}
    \definition{v.}{estampar (máquina de estamparia)}
  \synonymref{猛烈}{meng3lie4}
  \antonymref{缓}{huan3}
  \antonymref{慢}{man4}
  \antonymref{停}{ting2}
  \antonymref{止}{zhi3}
  \end{Phonetics}
\end{Entry}

\begin{Entry}{冲击}{6,5}{⼎,⼐}
  \begin{Phonetics}{冲击}{chong1ji1}[][HSK 6]
    \definition{v.}{chicotear; bater | correr; voar; atacar; assaltar; atacar bravamente em direção a um alvo predeterminado | chocar; metáfora para interferência ou golpe sério}
  \synonymref{冲刺}{chong1ci4}
  \synonymref{冲锋}{chong1feng1}
  \synonymref{出击}{chu1ji1}
  \synonymref{打击}{da3ji1}
  \synonymref{攻击}{gong1ji1}
  \synonymref{进攻}{jin4gong1}
  \synonymref{抨击}{peng1ji1}
  \synonymref{突击}{tu1ji1}
  \synonymref{袭击}{xi2ji1}
  \synonymref{撞击}{zhuang4ji1}
  \antonymref{后退}{hou4tui4}
  \antonymref{回避}{hui2bi4}
  \antonymref{稳定}{wen3ding4}
  \end{Phonetics}
\end{Entry}

\begin{Entry}{冲动}{6,6}{⼎,⼒}
  \begin{Phonetics}{冲动}{chong1dong4}[][HSK 5]
    \definition{adj.}{impulsivo; impetuoso}
    \definition{s.}{impulso; impetuosidade; impulso de movimento; fenômeno psicológico no qual as emoções são particularmente fortes e o controle racional é fraco}
    \definition{v.}{ficar animado; ser impetuoso; agir por impulso}
  \synonymref{冲击}{chong1ji1}
  \synonymref{刺激}{ci4ji1}
  \synonymref{感性}{gan3xing4}
  \synonymref{鼓动}{gu3dong4}
  \synonymref{激动}{ji1dong4}
  \synonymref{激发}{ji1fa1}
  \synonymref{激励}{ji1li4}
  \antonymref{遏制}{e4zhi4}
  \antonymref{静止}{jing4zhi3}
  \antonymref{冷静}{leng3jing4}
  \antonymref{理性}{li3xing4}
  \antonymref{理智}{li3zhi4}
  \antonymref{压制}{ya1zhi4}
  \antonymref{抑制}{yi4zhi4}
  \end{Phonetics}
\end{Entry}

\begin{Entry}{冲刺}{6,8}{⼎,⼑}
  \begin{Phonetics}{冲刺}{chong1ci4}[][HSK 7-9]
    \definition{v.}{arrancar; correr; disparar | metáfora para fazer o maior esforço ao se aproximar de uma meta ou estar prestes a ter sucesso}
  \end{Phonetics}
\end{Entry}

\begin{Entry}{冲洗}{6,9}{⼎,⽔}
  \begin{Phonetics}{冲洗}{chong1xi3}[][HSK 7-9]
    \definition{v.}{lavar; enxaguar | revelar filme; revelar e fixar o material fotossensível exposto}
  \end{Phonetics}
\end{Entry}

\begin{Entry}{冲突}{6,9}{⼎,⽳}
  \begin{Phonetics}{冲突}{chong1tu1}[][HSK 5]
    \definition{v.}{chocar"-se; entrar em conflito; conflitar | contradizer; duas coisas opostas que interferem uma na outra}
  \synonymref{抵触}{di3chu4}
  \synonymref{对立}{dui4li4}
  \synonymref{对峙}{dui4zhi4}
  \synonymref{矛盾}{mao2dun4}
  \synonymref{摩擦}{mo2ca1}
  \synonymref{争端}{zheng1duan1}
  \synonymref{争论}{zheng1lun4}
  \synonymref{争执}{zheng1zhi2}
  \antonymref{和解}{he2jie3}
  \antonymref{和谐}{he2xie2}
  \antonymref{平息}{ping2xi1}
  \antonymref{一致}{yi2zhi4}
  \end{Phonetics}
\end{Entry}

\begin{Entry}{冲浪}{6,10}{⼎,⽔}
  \begin{Phonetics}{冲浪}{chong1lang4}[][HSK 7-9]
    \definition{v.}{surfar; um esporte aquático em que os atletas surfam em pranchas especialmente construídas e deslizam ao longo das ondas | metáfora para navegar na \emph{Internet}}
  \end{Phonetics}
\end{Entry}

\begin{Entry}{冲破}{6,10}{⼎,⽯}
  \begin{Phonetics}{冲破}{chong1po4}
    \definition{v.}{romper; abrir uma brecha; estourar (abrir)}
  \seealsoref{突破}{tu1po4}
  \synonymref{打破}{da3po4}
  \antonymref{封锁}{feng1suo3}
  \end{Phonetics}
\end{Entry}

\begin{Entry}{冲锋}{6,12}{⼎,⾦}
  \begin{Phonetics}{冲锋}{chong1feng1}
    \definition{v.}{atacar; avançar; assaltar (frente de batalha); estar na vanguarda}
  \synonymref{冲刺}{chong1ci4}
  \synonymref{冲击}{chong1ji1}
  \antonymref{退却}{tui4que4}
  \end{Phonetics}
\end{Entry}

\begin{Entry}{冲撞}{6,15}{⼎,⼿}
  \begin{Phonetics}{冲撞}{chong1zhuang4}[][HSK 7-9]
    \definition{v.}{colidir; bater; sofrer impacto; voar; errar contra | ofender; ofender}
  \end{Phonetics}
\end{Entry}

%%%%%%%%%% 决 %%%%%%%%%%
\subsection*{决}\addcontentsline{loh}{figure}{决}

\begin{Entry}{决}{6}{⼎}
  \begin{Phonetics}{决}{jue2}
    \definition{v.}{decidir; determinar | executar uma pessoa | (de um dique, etc.) romper; desabar}
  \end{Phonetics}
\end{Entry}

\begin{Entry}{决不}{6,4}{⼎,⼀}
  \begin{Phonetics}{决不}{jue2bu4}[][HSK 5]
    \definition{adv.}{em hipótese alguma; nunca | definitivamente não; certamente não; sob nenhuma circunstância; de forma alguma}
  \seealsoref{决无}{jue2wu2}
  \synonymref{绝不}{jue2bu4}
  \antonymref{必定}{bi4ding4}
  \antonymref{一定}{yi2ding4}
  \end{Phonetics}
\end{Entry}

\begin{Entry}{决心}{6,4}{⼎,⼼}
  \begin{Phonetics}{决心}{jue2xin1}[][HSK 3]
    \definition{s.}{resolução; determinação; determinação inabalável}
    \definition{v.}{decidir"-se; tomar uma decisão; decidir fazer algo e não vacilar nem mudar de ideia}
  \seealsoref{决定}{jue2ding4}
  \synonymref{坚定}{jian1ding4}
  \synonymref{信念}{xin4nian4}
  \synonymref{信心}{xin4xin1}
  \synonymref{意志}{yi4zhi4}
  \antonymref{动摇}{dong4yao2}
  \antonymref{徘徊}{pai2huai2}
  \antonymref{犹豫}{you2yu4}
  \end{Phonetics}
\end{Entry}

\begin{Entry}{决无}{6,4}{⼎,⽆}
  \begin{Phonetics}{决无}{jue2wu2}
    \definition{adv.}{absolutamente não}
  \end{Phonetics}
\end{Entry}

\begin{Entry}{决议}{6,5}{⼎,⾔}
  \begin{Phonetics}{决议}{jue2yi4}[][HSK 7-9]
    \definition[项]{s.}{resolução | resultado}
  \end{Phonetics}
\end{Entry}

\begin{Entry}{决定}{6,8}{⼎,⼧}
  \begin{Phonetics}{决定}{jue2ding4}[][HSK 3]
    \definition{adj.}{decisivo; as leis objetivas levam as coisas a se desenvolverem e mudarem em determinada direção}
    \definition[项,个]{s.}{decisão; resolução; assuntos decididos}
    \definition{v.}{decidir; determinar; algo se torna a base ou o pré"-requisito para outra coisa; desempenha um papel dominante | decidir; resolver; tomar uma decisão; propor uma forma de agir}
  \seealsoref{决心}{jue2xin1}
  \synonymref{决策}{jue2ce4}
  \synonymref{决议}{jue2yi4}
  \synonymref{确定}{que4ding4}
  \antonymref{不定}{bu2ding4}
  \antonymref{迟疑}{chi2yi2}
  \antonymref{观望}{guan1wang4}
  \antonymref{无关}{wu2guan1}
  \antonymref{犹豫}{you2yu4}
  \end{Phonetics}
\end{Entry}

\begin{Entry}{决策}{6,12}{⼎,⽵}
  \begin{Phonetics}{决策}{jue2ce4}[][HSK 6]
    \definition{s.}{decisão política; decisão de importância estratégica; estratégia ou método de decisão}
    \definition{v.}{formular políticas; tomar uma decisão estratégica; decidir sobre uma estratégia ou abordagem}
  \synonymref{计划}{ji4hua4}
  \synonymref{决定}{jue2ding4}
  \synonymref{决议}{jue2yi4}
  \end{Phonetics}
\end{Entry}

\begin{Entry}{决赛}{6,14}{⼎,⾙}
  \begin{Phonetics}{决赛}{jue2sai4}[][HSK 3]
    \definition[场]{s.}{finais (de uma competição); em competições esportivas, a última partida ou rodada disputada para determinar a classificação}
  \antonymref{预赛}{yu4sai4}
  \end{Phonetics}
\end{Entry}

%%%%%%%%%% 况 %%%%%%%%%%
\subsection*{况}\addcontentsline{loh}{figure}{况}

\begin{Entry}{况}{7}{⼎}
  \begin{Phonetics}{况}{kuang4}
    \definition*{s.}{Sobrenome: Kuang}
    \definition{conj.}{além disso | mesmo; muito menos; sem mencionar}
    \definition{s.}{condição; situação}
    \definition{v.}{comparar}
  \end{Phonetics}
\end{Entry}

\begin{Entry}{况且}{7,5}{⼎,⼀}
  \begin{Phonetics}{况且}{kuang4qie3}
    \definition{conj.}{além disso; além do mais; orações de conexão para expressar uma relação progressiva}
  \end{Phonetics}
\end{Entry}

%%%%%%%%%% 冷 %%%%%%%%%%
\subsection*{冷}\addcontentsline{loh}{figure}{冷}

\begin{Entry}{冷}{7}{⼎}
  \begin{Phonetics}{冷}{leng3}[][HSK 1]
    \definition*{s.}{Sobrenome: Leng}
    \definition{adj.}{frio; baixa temperatura; sensação de frio | gelado; frio por natureza; sem entusiasmo; sem gentileza | desolado; pouco frequentado; quieto; sem agitação | negligenciado; indesejável; ignorado por todos | raro; estranho; incomum | feito em segredo; filmado de forma escondida; lançado secretamente}
    \definition{v.}{esfriar; resfriar | esfriar; congelar; tornar"-se indiferente, apático | ignorar}
  \synonymref{寒}{han2}
  \synonymref{寂}{ji4}
  \synonymref{凉}{liang2}
  \synonymref{清}{qing1}
  \antonymref{暖}{nuan3}
  \antonymref{热}{re4}
  \antonymref{炎}{yan2}
  \end{Phonetics}
\end{Entry}

\begin{Entry}{冷门}{7,3}{⼎,⾨}
  \begin{Phonetics}{冷门}{leng3men2}[][HSK 7-9]
    \definition{s.}{obscuro; impopular; originalmente, o termo se referia a um jogo de azar em que poucas pessoas apostavam; posteriormente, passou a ser usado metaforicamente para descrever algo que não chama a atenção ou que está fora de moda | um vencedor inesperado; azarão}
  \antonymref{热门}{re4men2}
  \end{Phonetics}
\end{Entry}

\begin{Entry}{冷气}{7,4}{⼎,⽓}
  \begin{Phonetics}{冷气}{leng3qi4}[][HSK 6]
    \definition[股,阵]{s.}{ar frio (ou fresco); correntes de ar frio | ar condicionado; ar resfriado por equipamento de refrigeração | ar condicionado; equipamentos de ar condicionado}
  \antonymref{暖气}{nuan3qi4}
  \end{Phonetics}
\end{Entry}

\begin{Entry}{冷水}{7,4}{⼎,⽔}
  \begin{Phonetics}{冷水}{leng3shui3}[][HSK 6]
    \definition[杯,瓶]{s.}{água fria | água não fervida}
  \synonymref{凉水}{liang2shui3}
  \antonymref{开水}{kai1shui3}
  \antonymref{热水}{re4shui3}
  \end{Phonetics}
\end{Entry}

\begin{Entry}{冷冻}{7,7}{⼎,⼎}
  \begin{Phonetics}{冷冻}{leng3dong4}[][HSK 7-9]
    \definition{v.}{congelar; refrigerar; congelar completamente}
  \end{Phonetics}
\end{Entry}

\begin{Entry}{冷战}{7,9}{⼎,⼽}
  \begin{Phonetics}{冷战}{leng3zhan4}[][HSK 7-9]
    \definition*{s.}{Guerra Fria; ações hostis fora de contexto de guerra, conduzidas internacionalmente}
    \definition{v.}{ignorar"-se completamente (duas pessoas que têm um bom relacionamento estão se ignorando, ou uma delas está ignorando a outra)}
  \end{Phonetics}
  \begin{Phonetics}{冷战}{leng3zhan5}
    \definition{s.}{Coloquial: arrepio; tremores ou calafrios repentinos devido ao frio ou ao medo}
  \seealsoref{冷颤}{leng3zhan5}
  \end{Phonetics}
\end{Entry}

\begin{Entry}{冷笑}{7,10}{⼎,⽵}
  \begin{Phonetics}{冷笑}{leng3xiao4}[][HSK 7-9]
    \definition{v.}{debochar; rir amargamente; sorrir com insatisfação, impotência, amargura, etc. | rir de forma sombria (com desdém); sorriso sardônico; uma risada que contém sarcasmo, insatisfação, impotência, desprezo, indiferença ou raiva | zombar}
  \end{Phonetics}
\end{Entry}

\begin{Entry}{冷淡}{7,11}{⼎,⽔}
  \begin{Phonetics}{冷淡}{leng3dan4}[][HSK 7-9]
    \definition{adj.}{indiferente; descreve a atitude de uma pessoa como desanimada ou indiferente | frouxo; sem florescimento; descreve uma loja ou estabelecimento similar como tendo pouco movimento, com pouquíssimas pessoas comprando coisas}
    \definition{v.}{ignorar alguém completamente; tratar alguém com hostilidade; fazer alguém se sentir indesejado}
  \end{Phonetics}
\end{Entry}

\begin{Entry}{冷落}{7,12}{⼎,⾋}
  \begin{Phonetics}{冷落}{leng3luo4}[][HSK 7-9]
    \definition{adj.}{desolado; deserto; pouco frequentado; silencioso; sem animação}
    \definition{v.}{ignorar alguém completamente; ser tratado com frieza}
  \end{Phonetics}
\end{Entry}

\begin{Entry}{冷漠}{7,13}{⼎,⽔}
  \begin{Phonetics}{冷漠}{leng3mo4}[][HSK 7-9]
    \definition{adj.}{desinteressado; indiferente; frio e distante; descreve alguém que não tem compaixão pelas desgraças alheias ou pelo sofrimento na sociedade; alguém que é indiferente e insensível (às pessoas ou às coisas)}
  \end{Phonetics}
\end{Entry}

\begin{Entry}{冷酷}{7,14}{⼎,⾣}
  \begin{Phonetics}{冷酷}{leng3ku4}[][HSK 7-9]
    \definition{adj.}{sombrio; insensível; cruel; (ao lidar com pessoas) frio e severo}
  \end{Phonetics}
\end{Entry}

\begin{Entry}{冷酷无情}{7,14,4,11}{⼎,⾣,⽆,⼼}
  \begin{Phonetics}{冷酷无情}{leng3ku4-wu2qing2}[][HSK 7-9]
    \definition{expr.}{de sangue frio; insensível; frio; com cruel frieza; sem afeto; de coração frio}
  \end{Phonetics}
\end{Entry}

\begin{Entry}{冷静}{7,14}{⼎,⾭}
  \begin{Phonetics}{冷静}{leng3jing4}[][HSK 4]
    \definition{adj.}{calmo; descreve uma pessoa que consegue ficar atenta em uma situação importante ou de emergência e não toma decisões aleatórias por causa de seus sentimentos no momento | (lugar) tranquilo; quieto; deserto}
  \synonymref{安静}{an1jing4}
  \synonymref{安宁}{an1ning2}
  \synonymref{从容}{cong2rong2}
  \synonymref{寂静}{ji4jing4}
  \synonymref{理智}{li3zhi4}
  \synonymref{宁静}{ning2jing4}
  \synonymref{平静}{ping2jing4}
  \antonymref{暴躁}{bao4zao4}
  \antonymref{冲动}{chong1dong4}
  \antonymref{疯狂}{feng1kuang2}
  \antonymref{慌张}{huang1zhang1}
  \antonymref{激动}{ji1dong4}
  \antonymref{焦虑}{jiao1lv4}
  \antonymref{热烈}{re4lie4}
  \antonymref{热闹}{re4nao5}
  \end{Phonetics}
\end{Entry}

\begin{Entry}{冷颤}{7,19}{⼎,⾴}
  \begin{Phonetics}{冷颤}{leng3zhan5}
    \definition{s.}{Coloquial: arrepio}
  \seealsoref{冷战}{leng3zhan5}
  \end{Phonetics}
\end{Entry}

%%%%%%%%%% 冻 %%%%%%%%%%
\subsection*{冻}\addcontentsline{loh}{figure}{冻}

\begin{Entry}{冻}{7}{⼎}
  \begin{Phonetics}{冻}{dong4}[][HSK 5]
    \definition*{s.}{Sobrenome: Dong}
    \definition{s.}{geleia; gelatina}
    \definition{v.}{congelar; ser congelado | ficar com frio ou sentir frio}
  \end{Phonetics}
\end{Entry}

\begin{Entry}{冻结}{7,9}{⼎,⽷}
  \begin{Phonetics}{冻结}{dong4jie2}[][HSK 7-9]
    \definition{adj.}{congelado}
    \definition{s.}{(salários, preços, etc.) congelamento; bloqueio de fluxo ou mudança}
    \definition{v.}{congelar}
  \end{Phonetics}
\end{Entry}

%%%%%%%%%% 净 %%%%%%%%%%
\subsection*{净}\addcontentsline{loh}{figure}{净}

\begin{Entry}{净}{8}{⼎}
  \begin{Phonetics}{净}{jing4}[][HSK 6]
    \definition{adj.}{limpo | (depois de um verbo) terminado; sem nada sobrando | líquido | vazio; oco; nu}
    \definition{adv.}{todo; o tempo todo | somente; meramente; nada além de | inteiramente; indica puro e nada mais}
    \definition{s.}{o ``rosto pintado'', comumente conhecido como Hualian (花脸), um tipo de personagem da ópera de Pequim, etc.}
    \definition{v.}{tornar limpo | limpar; lavar; esfregar para limpar}
  \seealsoref{花脸}{hua1lian3}
  \synonymref{纯}{chun2}
  \synonymref{光}{guang1}
  \synonymref{洁}{jie2}
  \synonymref{空}{kong1}
  \synonymref{空}{kong4}
  \synonymref{清}{qing1}
  \synonymref{总}{zong3}
  \antonymref{脏}{zang1}
  \end{Phonetics}
\end{Entry}

\begin{Entry}{净化}{8,4}{⼎,⼔}
  \begin{Phonetics}{净化}{jing4hua4}[][HSK 7-9]
    \definition{v.}{purificar; remover impurezas para purificar o objeto}
  \end{Phonetics}
\end{Entry}

%%%%%%%%%% 凄 %%%%%%%%%%
\subsection*{凄}\addcontentsline{loh}{figure}{凄}

\begin{Entry}{凄}{10}{⼎}
  \begin{Phonetics}{凄}{qi1}
    \definition{adj.}{frio; gelado | sombrio e desolado | triste; miserável; infeliz}
  \end{Phonetics}
\end{Entry}

\begin{Entry}{凄凉}{10,10}{⼎,⼎}
  \begin{Phonetics}{凄凉}{qi1liang2}[][HSK 7-9]
    \definition{adj.}{sombrio; desolador; solitário e desolado; miserável (frequentemente usado para descrever um ambiente ou cena) | miserável; sombrio e desolado; desolado e trágico}[夜晚的街道显得很凄凉。===As ruas estavam sombrias à noite.|他的眼神看起来很凄凉。===Seus olhos pareciam muito desolados.]
  \end{Phonetics}
\end{Entry}

\begin{Entry}{凄惨}{10,11}{⼎,⽕}
  \begin{Phonetics}{凄惨}{qi1can3}
    \definition{adj.}{trágico; miserável; deplorável; desolado e trágico}
  \seealsoref{悲惨}{bei1can3}
  \synonymref{惨痛}{can3tong4}
  \synonymref{凄凉}{qi1liang2}
  \antonymref{光辉}{guang1hui1}
  \antonymref{欢乐}{huan1le4}
  \antonymref{开心}{kai1//xin1}
  \antonymref{快乐}{kuai4le4}
  \antonymref{幸福}{xing4fu2}
  \antonymref{愉快}{yu2kuai4}
  \end{Phonetics}
\end{Entry}

%%%%%%%%%% 准 %%%%%%%%%%
\subsection*{准}\addcontentsline{loh}{figure}{准}

\begin{Entry}{准}{10}{⼎}
  \begin{Phonetics}{准}{zhun3}[][HSK 3]
    \definition{adj.}{exato; preciso; algo determinado a ser imutável | preciso; exato; correto | perto; parcialmente; quase; próximo de algo em termos de padrão}
    \definition{adv.}{definitivamente; certamente}
    \definition{pref.}{quasi-; para-}
    \definition{prep.}{de acordo com; baseado em}
    \definition{s.}{norma; padrão; critério | confiança certa; uma ideia definida, certeza, etc. (geralmente usada depois de 有 ou 没有)}
    \definition{v.}{autorizar; conceder; consentir; permitir}
  \seealsoref{没有}{mei2you5}
  \seealsoref{同意}{tong2yi4}
  \seealsoref{有}{you3}
  \synonymref{对}{dui4}
  \synonymref{许}{xu3}
  \synonymref{允}{yun3}
  \antonymref{错}{cuo4}
  \antonymref{误}{wu4}
  \end{Phonetics}
\end{Entry}

\begin{Entry}{准则}{10,6}{⼎,⼑}
  \begin{Phonetics}{准则}{zhun3ze2}[][HSK 7-9]
    \definition[条]{s.}{norma; padrão; critério; regra; fórmula; diretriz; os princípios em que se baseiam a fala e as ações}
  \synonymref{标准}{biao1zhun3}
  \synonymref{法规}{fa3gui1}
  \synonymref{法则}{fa3ze2}
  \synonymref{规则}{gui1ze2}
  \synonymref{规矩}{gui1ju5}
  \synonymref{原则}{yuan2ze2}
  \end{Phonetics}
\end{Entry}

\begin{Entry}{准许}{10,6}{⼎,⾔}
  \begin{Phonetics}{准许}{zhun3xu3}[][HSK 7-9]
    \definition{v.}{permitir; autorizar; concordar com os pedidos de outras pessoas}
  \synonymref{承诺}{cheng2nuo4}
  \synonymref{答应}{da1ying5}
  \synonymref{批准}{pi1zhun3}
  \synonymref{容许}{rong2xu3}
  \synonymref{同意}{tong2yi4}
  \synonymref{许可}{xu3ke3}
  \synonymref{愿意}{yuan4yi5}
  \synonymref{允许}{yun3xu3}
  \antonymref{不许}{bu4xu3}
  \antonymref{禁止}{jin4zhi3}
  \end{Phonetics}
\end{Entry}

\begin{Entry}{准时}{10,7}{⼎,⽇}
  \begin{Phonetics}{准时}{zhun3shi2}[][HSK 4]
    \definition{adj.}{pontual}
    \definition{adv.}{na hora certa; dentro do prazo; no horário especificado}
  \synonymref{按期}{an4qi1}
  \synonymref{按时}{an4shi2}
  \synonymref{定时}{ding4shi2}
  \synonymref{及时}{ji2shi2}
  \synonymref{精确}{jing1que4}
  \synonymref{正确}{zheng4que4}
  \synonymref{准确}{zhun3que4}
  \antonymref{迟到}{chi2//dao4}
  \antonymref{耽搁}{dan1ge5}
  \antonymref{误点}{wu4//dian3}
  \antonymref{延误}{yan2wu4}
  \end{Phonetics}
\end{Entry}

\begin{Entry}{准备}{10,8}{⼎,⼡}
  \begin{Phonetics}{准备}{zhun3bei4}[][HSK 1]
    \definition{v.}{preparar; ficar pronto; planejar ou organizar com antecedência | pretender; planejar}
  \seealsoref{筹备}{chou2bei4}
  \synonymref{安排}{an1pai2}
  \synonymref{筹办}{chou2ban4}
  \synonymref{筹划}{chou2hua4}
  \synonymref{打算}{da3suan5}
  \synonymref{计划}{ji4hua4}
  \synonymref{盘算}{pan2suan5}
  \synonymref{预备}{yu4bei4}
  \antonymref{放弃}{fang4qi4}
  \end{Phonetics}
\end{Entry}

\begin{Entry}{准确}{10,12}{⼎,⽯}
  \begin{Phonetics}{准确}{zhun3que4}[][HSK 2]
    \definition{adj.}{exato; preciso; acurado; os resultados da ação são completamente consistentes com os resultados reais ou esperados}
  \synonymref{精确}{jing1que4}
  \synonymref{切实}{qie4shi2}
  \synonymref{确切}{que4qie4}
  \synonymref{确实}{que4shi2}
  \synonymref{确凿}{que4zao2}
  \synonymref{正确}{zheng4que4}
  \antonymref{错误}{cuo4wu4}
  \antonymref{大约}{da4yue1}
  \antonymref{大致}{da4zhi4}
  \antonymref{偏差}{pian1cha1}
  \end{Phonetics}
\end{Entry}

%%%%%%%%%% 凉 %%%%%%%%%%
\subsection*{凉}\addcontentsline{loh}{figure}{凉}

\begin{Entry}{凉}{10}{⼎}
  \begin{Phonetics}{凉}{liang2}[][HSK 2]
    \definition{adj.}{frio; gelado; ligeiramente fria (menos do que 冷) | sombrio; desolado; sem animação | desanimado; desapontado | usado para prevenir o calor e manter a temperatura amena; para proteção contra o calor}
    \definition{s.}{frio; refere"-se a um ambiente fresco ou a uma brisa fresca}
  \seealsoref{冷}{leng3}
  \synonymref{悲}{bei1}
  \synonymref{愁}{chou2}
  \synonymref{苦}{ku3}
  \synonymref{伤}{shang1}
  \antonymref{暖}{nuan3}
  \antonymref{热}{re4}
  \antonymref{温}{wen1}
  \antonymref{炎}{yan2}
  \end{Phonetics}
  \begin{Phonetics}{凉}{liang4}
    \definition{v.}{deixar algo esfriar; deixar um objeto quente descansar por um tempo para que a temperatura diminua}
  \synonymref{冷}{leng3}
  \antonymref{热}{re4}
  \antonymref{炎}{yan2}
  \end{Phonetics}
\end{Entry}

\begin{Entry}{凉水}{10,4}{⼎,⽔}
  \begin{Phonetics}{凉水}{liang2shui3}[][HSK 3]
    \definition[杯]{s.}{água fria; água não aquecida | água não fervida}
  \synonymref{冷水}{leng3shui3}
  \antonymref{开水}{kai1shui3}
  \antonymref{热水}{re4shui3}
  \end{Phonetics}
\end{Entry}

\begin{Entry}{凉快}{10,7}{⼎,⼼}
  \begin{Phonetics}{凉快}{liang2kuai5}[][HSK 2]
    \definition{adj.}{fresco; refrescante}
    \definition{v.}{refrescar; refrescar"-se; deixar o corpo fresco e revigorado}
  \seealsoref{凉爽}{liang2shuang3}
  \synonymref{降温}{jiang4wen1}
  \synonymref{清凉}{qing1liang2}
  \synonymref{清爽}{qing1shuang3}
  \antonymref{闷热}{men1re4}
  \antonymref{暖和}{nuan3huo5}
  \antonymref{炎热}{yan2re4}
  \end{Phonetics}
\end{Entry}

\begin{Entry}{凉爽}{10,11}{⼎,⽘}
  \begin{Phonetics}{凉爽}{liang2shuang3}[][HSK 7-9]
    \definition{adj.}{agradável e fresco; agradavelmente fresco; não está quente, é uma sensação agradável}
  \end{Phonetics}
\end{Entry}

\begin{Entry}{凉鞋}{10,15}{⼎,⾰}
  \begin{Phonetics}{凉鞋}{liang2xie2}[][HSK 6]
    \definition[双,只]{s.}{sandália; alpargata; alpercata; alparca ; sapatos de verão com cabedal ventilado}
  \end{Phonetics}
\end{Entry}

%%%%%%%%%% 凌 %%%%%%%%%%
\subsection*{凌}\addcontentsline{loh}{figure}{凌}

\begin{Entry}{凌}{10}{⼎}
  \begin{Phonetics}{凌}{ling2}
    \definition*{s.}{Sobrenome: Ling}
    \definition{s.}{gelo}
    \definition{v.}{insultar; invadir; violar; abusar | elevar"-se alto; subir; ascender | abordar; aproximar"-se}
  \end{Phonetics}
\end{Entry}

\begin{Entry}{凌晨}{10,11}{⼎,⽇}
  \begin{Phonetics}{凌晨}{ling2chen2}[][HSK 7-9]
    \definition[个]{s.}{antes do amanhecer; nas primeiras horas da manhã; com a aproximação do amanhecer}
  \end{Phonetics}
\end{Entry}

%%%%%%%%%% 减 %%%%%%%%%%
\subsection*{减}\addcontentsline{loh}{figure}{减}

\begin{Entry}{减}{11}{⼎}
  \begin{Phonetics}{减}{jian3}[][HSK 4]
    \definition*{s.}{Sobrenome: Jian}
    \definition{v.}{subtrair; remover uma parte da quantidade original | reduzir; diminuir; cortar}
  \synonymref{除}{chu2}
  \antonymref{加}{jia1}
  \antonymref{添}{tian1}
  \antonymref{增}{zeng1}
  \end{Phonetics}
\end{Entry}

\begin{Entry}{减少}{11,4}{⼎,⼩}
  \begin{Phonetics}{减少}{jian3shao3}[][HSK 4]
    \definition{v.}{cair; reduzir; diminuir; subtrair uma parte}
  \synonymref{减轻}{jian3qing1}
  \synonymref{降低}{jiang4di1}
  \synonymref{省略}{sheng3lve4}
  \synonymref{收缩}{shou1suo1}
  \synonymref{缩短}{suo1duan3}
  \synonymref{缩小}{suo1xiao3}
  \synonymref{削弱}{xue1ruo4}
  \antonymref{加强}{jia1qiang2}
  \antonymref{增多}{zeng1duo1}
  \antonymref{增加}{zeng1jia1}
  \antonymref{增强}{zeng1qiang2}
  \antonymref{增添}{zeng1tian1}
  \antonymref{增长}{zeng1zhang3}
  \end{Phonetics}
\end{Entry}

\begin{Entry}{减压}{11,6}{⼎,⼚}
  \begin{Phonetics}{减压}{jian3ya1}[][HSK 7-9]
    \definition{v.}{reduzir a pressão; descomprimir | reduzir o fardo | reduzir a pressão; despressurizar; descomprimir | relaxar}
  \end{Phonetics}
\end{Entry}

\begin{Entry}{减免}{11,7}{⼎,⼉}
  \begin{Phonetics}{减免}{jian3mian3}[][HSK 7-9]
    \definition{v.}{mitigar ou anular (uma punição) | reduzir ou isentar (impostos, etc.)}
  \end{Phonetics}
\end{Entry}

\begin{Entry}{减肥}{11,8}{⼎,⾁}
  \begin{Phonetics}{减肥}{jian3//fei2}[][HSK 4]
    \definition{v.+compl.}{perder peso; dieta, exercícios, medicamentos, massagem, cirurgia, etc., para reduzir o excesso de gordura corporal, de modo que o grau de obesidade seja reduzido}
  \end{Phonetics}
\end{Entry}

\begin{Entry}{减削}{11,9}{⼎,⼑}
  \begin{Phonetics}{减削}{jian3xue1}
    \definition{v.}{reduzir; cortar; reduzir e alterar}
  \end{Phonetics}
\end{Entry}

\begin{Entry}{减轻}{11,9}{⼎,⾞}
  \begin{Phonetics}{减轻}{jian3qing1}[][HSK 5]
    \definition{v.}{aliviar; remeter; clarear; facilitar; mitigar}
  \seealsoref{减弱}{jian3ruo4}
  \synonymref{淡化}{dan4hua4}
  \synonymref{减少}{jian3shao3}
  \synonymref{降低}{jiang4di1}
  \synonymref{削弱}{xue1ruo4}
  \antonymref{加剧}{jia1ju4}
  \antonymref{加重}{jia1zhong4}
  \antonymref{增多}{zeng1duo1}
  \antonymref{增加}{zeng1jia1}
  \antonymref{增添}{zeng1tian1}
  \end{Phonetics}
\end{Entry}

\begin{Entry}{减弱}{11,10}{⼎,⼸}
  \begin{Phonetics}{减弱}{jian3ruo4}[][HSK 7-9]
    \definition{v.}{reduzir | enfraquecer}
  \end{Phonetics}
\end{Entry}

\begin{Entry}{减速}{11,10}{⼎,⾡}
  \begin{Phonetics}{减速}{jian3//su4}[][HSK 7-9]
    \definition{v.+compl.}{diminuir a velocidade; desacelerar; retardar | moderar; reduzir a velocidade; reduzir a marcha; atrasar; desacelerar; retardar}
  \antonymref{加速}{jia1su4}
  \end{Phonetics}
\end{Entry}

%%%%%%%%%% 凑 %%%%%%%%%%
\subsection*{凑}\addcontentsline{loh}{figure}{凑}

\begin{Entry}{凑}{11}{⼎}
  \begin{Phonetics}{凑}{cou4}[][HSK 7-9]
    \definition{v.}{reunir; coletar; ajuntar | acontecer por acaso; aproveitar; esbarrar em; alcançar; tirar vantagem de | aproximar; mover"-se para perto de}
  \end{Phonetics}
\end{Entry}

\begin{Entry}{凑巧}{11,5}{⼎,⼯}
  \begin{Phonetics}{凑巧}{cou4qiao3}[][HSK 7-9]
    \definition{adj.}{afortunado; sortudo; coincidente; significa que é o momento certo ou que algo que você quer ou não quer está acontecendo}
  \end{Phonetics}
\end{Entry}

\begin{Entry}{凑合}{11,6}{⼎,⼝}
  \begin{Phonetics}{凑合}{cou4he5}[][HSK 7-9]
    \definition{v.}{contentar"-se com algo; ser razoável; ser razoavelmente bom, mas não excelente; aceitar relutantemente coisas ou condições de um nível ou grau inferior | improvisar | reunir}
  \end{Phonetics}
\end{Entry}

%%%%%%%%%% 凝 %%%%%%%%%%
\subsection*{凝}\addcontentsline{loh}{figure}{凝}

\begin{Entry}{凝}{16}{⼎}
  \begin{Phonetics}{凝}{ning2}
    \definition{adv.}{atentamente; atenção fixa}
    \definition{v.}{congelar; coagular; coalhar; condensar | contemplar; olhar pensativamente}
  \end{Phonetics}
\end{Entry}

\begin{Entry}{凝固}{16,8}{⼎,⼞}
  \begin{Phonetics}{凝固}{ning2gu4}[][HSK 7-9]
    \definition{v.}{estagnar; parar, não se move mais nem muda de direção}
  \end{Phonetics}
\end{Entry}

\begin{Entry}{凝聚}{16,14}{⼎,⽿}
  \begin{Phonetics}{凝聚}{ning2ju4}[][HSK 7-9]
    \definition{v.}{condensar (o vapor); coagular; coalhar (os fluidos) | reunir; acumular; juntar}
  \end{Phonetics}
\end{Entry}

%%%%% EOF %%%%%

